\documentclass[11pt]{article}
\usepackage[margin=1in]{geometry}
\usepackage{parskip}
\usepackage{enumitem}
\usepackage[T1]{fontenc}
\usepackage{booktabs}

\title{\textbf{Assessment of PROVSAFE}}
\author{}
\date{}

\begin{document}
\maketitle

\section*{Motivation}

Tool-using LLM agents increasingly interact with email, browsers, APIs, databases, and external documents. In these settings, the agent does not reason only over the user prompt; it also reasons over tool outputs that may contain attacker-controlled text. This creates the core problem of indirect prompt injection: malicious instructions can be embedded in retrieved content, web pages, emails, or files and then influence downstream tool calls.

The paper identifies an important gap in existing defenses. Many current approaches focus on keyword filtering, prompt rewriting, secondary LLM moderation, or action validation without explicitly modeling where the triggering information came from. As a result, they may fail to distinguish legitimate user instructions from malicious content introduced through tools.

\section*{Research Insight}

 \textbf{Authorization decisions for LLM agent actions should depend not only on the requested action, but also on the provenance of the information that produced the action arguments.}

This is a valuable insight because it changes the security decision from a purely content-based or pattern-based check into a lineage-aware enforcement problem. Instead of asking only whether a tool call appears dangerous, the system asks whether the arguments for that call originated from trusted user intent, untrusted external data, or a mixture of both.

A second insight is that provenance in LLM agents cannot be handled exactly like classical taint tracking. Since LLMs may paraphrase, summarize, or partially transform content before generating a tool call, provenance resolution must tolerate semantic variation while remaining conservative in enforcement. This gives the work more depth than a simple wrapper around existing policy systems.

\section*{Scientific Contribution}

The strongest contribution is not merely that it implements another defense framework, but that it introduces a \textbf{security mechanism} for tool-using LLM agents based on provenance-gated authorization.

The contribution can be understood in several parts:
\begin{itemize}[leftmargin=1.5em]
    \item a \textbf{conceptual contribution}, by arguing that provenance is a missing security primitive for LLM agents,
    \item a \textbf{methodological contribution}, by formulating provenance-aware policy enforcement at the tool boundary with formal properties (Complete Mediation, Policy Determinism, Conservative Trust Propagation),
    \item a \textbf{systems contribution}, by implementing a working framework that tracks provenance via a W3C PROV-DM DAG with a two-element trust lattice and three-stage resolution pipeline (substring matching, embedding similarity, conservative default),
    \item and an \textbf{empirical contribution}, by evaluating the approach across 25,000+ trials (5 models spanning five vendor families, 5 systems including Taint-Everything, 5 reps) with external validation on the InjecAgent benchmark (1,054 cases).
\end{itemize}

This is stronger than a paper whose novelty lies mainly in integration or benchmarking. Here, the paper has a central mechanism that is both technically meaningful and conceptually distinct.

\section*{Method}

The paper defines a provenance-aware architecture in which agent inputs, tool outputs, intermediate transformations, and final tool-call arguments are linked through a provenance graph. Trust labels are propagated through this structure via a lattice meet operator, and enforcement rules are evaluated at the moment a tool call is about to be executed.

This design is strong for several reasons. First, it places the defense at the \textbf{tool-call boundary}, which is a practical enforcement point because it does not require changing the internal weights or decoding process of the model. Second, it combines provenance tracking with declarative YAML policy checks, which makes the system more auditable than purely model-based defenses. Third, it includes embedding-based provenance resolution (all-MiniLM-L6-v2) for paraphrased content, which addresses a realistic problem in LLM-generated outputs and achieves 100\% provenance recall across 96 paraphrase tests.

\section*{Key Results}

\begin{table}[h]
\centering
\small
\begin{tabular}{lrrr}
\toprule
\textbf{System} & \textbf{ASR-IA $\downarrow$} & \textbf{TSR $\uparrow$} & \textbf{FPR $\downarrow$} \\
\midrule
No Defense        & 24.56\% & 98.38\% & 1.62\% \\
Pattern Filter    & 7.33\%  & 99.12\% & 0.88\% \\
Policy-Only       & 1.72\%  & 96.75\% & 3.25\% \\
Taint-Everything  & 0.00\%  & 9.00\% & 91.00\% \\
\textbf{PROVSAFE} & \textbf{1.02\%} & \textbf{95.00\%} & \textbf{5.00\%} \\
\bottomrule
\end{tabular}
\caption{primary benchmark: 26,500 trials (212 scenarios $\times$ 5 models $\times$ 5 systems $\times$ 5 reps). PROVSAFE achieves a 24$\times$ reduction in attack success versus No Defense while maintaining 95\% usability---unlike Taint-Everything, which achieves 0.00\% ASR at the cost of 9\% TSR (unusable).}
\end{table}

\begin{itemize}[leftmargin=1.5em]
    \item \textbf{24$\times$ reduction} in attack success rate versus undefended baselines (24.56\% $\to$ 1.02\%).
    \item \textbf{Taint-Everything is secure but unusable}: achieves 0.00\% ASR but collapses TSR to 9.00\%---PROVSAFE matches its security (1.02\% vs.\ 0.00\%) while preserving 95\% usability.
    \item \textbf{Provenance halves residual risk}: Policy-Only achieves 1.72\% ASR; adding provenance reduces this to 1.02\%---the attacks that provenance catches are precisely those where syntax-based rules cannot distinguish attacker-injected from user-initiated arguments.
    \item \textbf{0\% ASR on 5 of 9 attack categories}: direct injection, device-state injection, file-content injection, role-play jailbreak, and NL-argument injection are fully blocked.
    \item \textbf{External validation}: On the InjecAgent benchmark (1,054 independent cases, 17 user tools, 63 attacker tools), PROVSAFE achieves 4.12\% ASR versus 20.83\% No Defense---a 5$\times$ reduction on unseen tools without any policy modification.
    \item \textbf{Paraphrase robustness}: Embedding-enhanced provenance achieves 100\% recall on 96 paraphrase tests (vs.\ 40.6\% for substring matching alone), with $<$5\,ms overhead.
    \item \textbf{Adaptive adversary}: 14/15 white-box attacks blocked (6.7\% adaptive ASR). The single bypass is a policy gap (MEDIUM-risk exfiltration channel), not a provenance failure.
    \item \textbf{Consistent across models}: ASR ranges 0.00\%--1.74\% across five models spanning five vendor families (Meta, Alibaba, Google, Microsoft, OpenAI), with \textbf{0.35\% ASR-IA on GPT-4o-mini} (99.5\% TSR).
    \item \textbf{Concrete proof}: On InjecAgent case \texttt{dh\_base\_0016}, an injected web page tricks the LLM into calling \texttt{AugustSmartLockGrantGuestAccess}---a legitimate tool. Policy-Only has no rule to block it. PROVSAFE traces the arguments to the untrusted web response and denies the call. This is the gap provenance fills.
\end{itemize}

\section*{Security--Usability Tradeoff}

The Taint-Everything baseline demonstrates why blanket tainting is not a viable alternative: it achieves 0.00\% ASR but at the cost of 9.00\% TSR---a 91\% false positive rate that renders the system unusable. PROVSAFE achieves nearly the same security while preserving 95\% usability. Compared to Policy-Only, PROVSAFE adds $\sim$2\% more false positives (5.00\% vs.\ 3.25\%) to reduce the residual attack rate (1.02\% vs.\ 1.72\%). All false positives trace to specific, auditable policy rules---not provenance misclassification---meaning they can be reduced through policy refinement without weakening security.

\section*{Target Venue}

\textbf{IEEE Transactions on Dependable and Secure Computing (TDSC)} is the primary target:
\begin{itemize}[leftmargin=1.5em]
    \item Journal scope matches exactly: provenance + LLM security + formal properties.
    \item The formal framework (W3C PROV-DM, trust lattice, proof sketches for 3 theorems) fits TDSC's expectations for rigor.
    \item 26,500+-trial evaluation (5 models across five vendor families, 5 systems including Taint-Everything) with external benchmark validation exceeds typical TDSC empirical standards.
    \item Journal format allows revision rounds to address reviewer concerns (larger models, user study).
\end{itemize}

\section*{Pitch}

\begin{quote}
\textit{PROVSAFE introduces provenance-aware authorization at the tool-call boundary as a security primitive for defending tool-using LLM agents against indirect prompt injection. Across 26,500+ trials (five models, five vendor families) and 1,054 external benchmark cases, provenance-gated enforcement reduces attack success to 1.02\% while maintaining 95\% task success---achieving near-identical security to Taint-Everything (0.00\% ASR) without its catastrophic usability collapse (9.00\% TSR). On InjecAgent, policy enforcement without provenance fails at 22.67\% ASR (no better than no defense), while PROVSAFE reduces ASR to 4.12\%, proving that tracking where arguments came from is the missing piece in LLM agent security.}
\end{quote}

\end{document}
